\section{Requisitos}
\label{sec:requisitos}

As Tabelas~\ref{tab:rf} e~\ref{tab:rnf} especificam requisitos-alvo do
artefato completo; sua presença não significa que todos estejam implementados
na versão de andamento. Os atores preliminares são: \textbf{frentista}, que solicita autorização;
\textbf{supervisor}, que trata exceções justificadas; \textbf{administrador},
que mantém usuários, identificadores dos conjuntos de bicos e política;
\textbf{auditor/pesquisador}, que acessa
dados minimizados; \textbf{provedor veicular}, simulado ou autorizado; e
\textbf{controlador de bancada}, que executa comandos restritos. Cliente e
titular de dados não são usuários do protótipo, mas são partes interessadas em
segurança, transparência e proteção de dados.

\begin{table}[H]
  \centering
  \caption{Requisitos funcionais preliminares}
  \label{tab:rf}
  \small
  \begin{tabularx}{\textwidth}{@{}L{1.1cm}Y@{}}
    \toprule
    \textbf{ID} & \textbf{Requisito} \\
    \midrule
    RF01 & Receber identificador idempotente da solicitação, placa, combustível, frentista e conjunto de bicos (\path{pump_id}), validando os campos. \\
    RF02 & Consultar um \textit{VehicleDataProvider} sem expor credenciais ao cliente web. \\
    RF03 & Converter códigos externos em categorias internas versionadas e registrar a versão da regra. \\
    RF04 & Retornar \path{AUTHORIZED}, \path{DENIED} ou \path{REVIEW_REQUIRED} com motivo padronizado. \\
    RF05 & Autorizar temporariamente a abertura somente da trava do suporte associado ao produto solicitado e compatível. \\
    RF06 & Desenergizar todas as saídas de abertura após expiração, encerramento, reinicialização ou perda de comunicação. \\
    RF07 & Permitir override somente a supervisor autenticado, com justificativa e validade limitada. \\
    RF08 & Rejeitar comando duplicado ou vencido e registrar a confirmação do dispositivo. \\
    RF09 & Manter trilha de auditoria minimizada, sem placa ou resposta veicular bruta no log padrão. \\
    RF10 & Operar com provedores e gateways simulados para desenvolvimento e demonstração sem dado real. \\
    \bottomrule
  \end{tabularx}
  \\[2pt]\fontepropria
\end{table}

\begin{table}[H]
  \centering
  \caption{Requisitos não funcionais e restrições preliminares}
  \label{tab:rnf}
  \small
  \begin{tabularx}{\textwidth}{@{}L{1.2cm}L{2.9cm}Y@{}}
    \toprule
    \textbf{ID} & \textbf{Atributo} & \textbf{Critério} \\
    \midrule
    RNF01 & Segurança funcional & Estado inicial e estado após falha devem ser \path{LOCK_ALL}. \\
    RNF02 & Privacidade & Usar minimização e não persistir placa, chassi, Renavam, proprietário ou resposta bruta por padrão. \\
    RNF03 & Segurança lógica & Credenciais externas ao repositório, autenticação individual e menor privilégio. \\
    RNF04 & Auditabilidade & Toda decisão e override deve registrar horário, ator, motivo, versão da política e resultado do dispositivo. \\
    RNF05 & Interoperabilidade & Provedor e gateway devem respeitar interfaces substituíveis, sem acoplamento do domínio à SENATRAN ou ao ESP32. \\
    RNF06 & Desempenho & Medir mediana, P95 e máximo; o limite operacional será definido após linha de base, sem número arbitrário nesta versão. \\
    RNF07 & Reprodutibilidade & Casos sintéticos, configuração de bancada e testes devem ser versionados. \\
    RNF08 & Segurança física & Operar apenas em bancada seca e de baixa tensão; não comandar bombeamento nem automação de pista; instalação real fica fora do escopo. \\
    \bottomrule
  \end{tabularx}
  \\[2pt]\fontepropria
\end{table}

Na versão verificada em 31 de julho de 2026, o backend já executa decisões com
provedor e intertravamento simulados, rejeita comandos duplicados ou vencidos e
mantém auditoria minimizada. A repetição de uma solicitação com o mesmo
\path{client_request_id} e o mesmo conteúdo recupera a decisão existente sem
novo acionamento; conteúdo divergente gera conflito. Entretanto, RF03, RF07,
RF08, RNF03 e RNF04 estão
apenas parcialmente atendidos: ainda não há versão de política persistida,
autenticação individual nem registro integrado de confirmação física. O
override usa um token Bearer compartilhado e recebe o identificador de
supervisor informado na requisição; isso demonstra o fluxo, mas não comprova a
identidade do ator. O firmware experimental emite confirmações, e o backend já
possui um adaptador MQTT que valida identificador, sequência, sessão, estado e
bico antes de aceitar o ACK. A integração ao broker ainda não foi ensaiada ao
vivo, e o ACK não constitui uma entidade persistida independente. Mesmo quando
coerente, ele confirma o comando e o estado lógico da saída GPIO, não a posição
mecânica da trava ou do bico.

Os requisitos serão validados com o orientador e, depois das autorizações
necessárias, com representantes dos papéis operacionais. Alterações serão
rastreáveis até a motivação, o cenário e o teste correspondente.
